\subsection{Kwalifikacja metodą pogłosową: układ i opracowanie}
Jako niezależne porównanie z oszacowaniem $A=\alpha S_v$ wykorzystano
czasy pogłosu z 12 plików w katalogu
\texttt{data/raw/grupa 1 Wibro tech 11 30/}.
Zgodnie z przebiegiem ćwiczenia zastosowano dwa ustawienia mikrofonu.
Dla każdego z nich głośnik ustawiano kolejno w sześciu położeniach
po przeciwnej stronie pomieszczenia; następnie zamieniono strony.
Serię plików \texttt{z1m1--z1m6} przypisano do pierwszego ustawienia,
a \texttt{z2m1--z2m6} do drugiego, zgodnie z opisem układu pomiarowego.
Nazwy plików zachowano, a oznaczenia $M_1$, $M_2$ i $G_1$--$G_6$
na schemacie służą opisowi fizycznych położeń.
\begin{figure}[H]
\centering
\begin{tikzpicture}[scale=0.93,every node/.style={font=\scriptsize}]
\begin{scope}
\draw[thick] (0,0) rectangle (6,3.75);
\node at (3,4.15) {Ustawienie 1};
\fill[accent] (0.8,1.88) circle (0.10);
\node[below] at (0.8,1.7) {$M_1$};
\foreach \x/\y/\n in {3.7/0.8/1,4.9/0.8/2,3.7/1.88/3,4.9/1.88/4,3.7/2.95/5,4.9/2.95/6}{
  \draw[fill=gray!20] (\x-0.12,\y-0.12) rectangle (\x+0.12,\y+0.12);
  \node[above] at (\x,\y+0.15) {$G_{\n}$};
}
\draw[<->] (0,-0.3)--(6,-0.3) node[midway,below] {8 m};
\node[rotate=90] at (-0.35,1.88) {5 m};
\end{scope}
\begin{scope}[xshift=7.1cm]
\draw[thick] (0,0) rectangle (6,3.75);
\node at (3,4.15) {Ustawienie 2};
\fill[accent] (5.2,1.88) circle (0.10);
\node[below] at (5.2,1.7) {$M_2$};
\foreach \x/\y/\n in {1.1/0.8/1,2.3/0.8/2,1.1/1.88/3,2.3/1.88/4,1.1/2.95/5,2.3/2.95/6}{
  \draw[fill=gray!20] (\x-0.12,\y-0.12) rectangle (\x+0.12,\y+0.12);
  \node[above] at (\x,\y+0.15) {$G_{\n}$};
}
\draw[<->] (0,-0.3)--(6,-0.3) node[midway,below] {8 m};
\end{scope}
\end{tikzpicture}
\caption{Dwa ustawienia pomiaru pogłosu w rzucie z góry.
Koło oznacza mikrofon, kwadraty sześć kolejnych położeń głośnika
(nie sześć źródeł pracujących jednocześnie). Po zmianie ustawienia
zamieniono strony pomieszczenia. Układ punktów jest poglądowy:
rysunek nie określa zmierzonych odległości ani dokładnej kolejności ruchu.}
\label{fig:poglos-ustawienia}
\end{figure}


Eksporty zawierają EDT, T10, T15, T20, T30 i RT oraz współczynniki
korelacji regresji. Do porównania przyjęto jednolicie parametr T20,
aby nie mieszać różnych zakresów regresji pomiędzy częstotliwościami.
T20 jest estymacją czasu zaniku o 60 dB na podstawie nachylenia
w zakresie od $-5$ do $-25$ dB; wartości z eksportu nie mnoży się
ponownie przez trzy \autocite{bkReverb}.
Czasy uśredniono arytmetycznie po 12 kombinacjach położeń:
\begin{equation}
\overline T(f)=\frac{1}{12}\sum_{j=1}^{2}\sum_{i=1}^{6}T_{20,ji}(f).
\end{equation}
To średnia przestrzenna, a nie uśrednianie energetyczne poziomów dźwięku.
Sam eksport nie zawiera przebiegów zaniku ani pełnych informacji
o wzbudzeniu i torze pomiarowym, dlatego nie odtwarzano ich z domysłu.

\subsection{Czasy pogłosu i chłonność}
Pliki zawierają pasma tercjowe od 25 Hz do 20 kHz.
W tabeli pokazano wybrane pasma o częstotliwościach środkowych
125--8000 Hz; nie są to wyniki analizy oktawowej.
Nie uśredniano czasów trzech tercji w celu zastąpienia pomiaru oktawowego.
Komplet wartości, w tym braki oznaczone w eksporcie przez \texttt{--},
zachowano w przetworzonym zestawieniu danych.
\begin{table}[H]
\centering\small
\caption{T20 w sekundach: wybrane pasma tercjowe. Nazwy wierszy odpowiadają plikom.}
\begin{tabular}{@{}l r r r r r r r@{}}
\toprule
Plik & 125 Hz & 250 Hz & 500 Hz & 1 kHz & 2 kHz & 4 kHz & 8 kHz\\
\midrule
z1m1 & \num{0.630} & \num{0.936} & \num{0.713} & \num{0.671} & \num{0.681} & \num{0.673} & \num{0.547} \\
z1m2 & \num{0.546} & \num{0.848} & \num{0.746} & \num{0.716} & \num{0.799} & \num{0.606} & \num{0.495} \\
z1m3 & \num{0.645} & \num{0.605} & \num{0.726} & \num{0.758} & \num{0.612} & \num{0.606} & \num{0.485} \\
z1m4 & \num{0.447} & \num{0.913} & \num{0.712} & \num{0.755} & \num{0.713} & \num{0.666} & \num{0.555} \\
z1m5 & \num{0.516} & \num{0.935} & \num{0.662} & \num{0.737} & \num{0.635} & \num{0.692} & \num{0.563} \\
z1m6 & \num{0.480} & \num{0.731} & \num{0.805} & \num{0.874} & \num{0.651} & \num{0.622} & \num{0.458} \\
z2m1 & \num{0.706} & \num{0.622} & \num{0.639} & \num{0.819} & \num{0.727} & \num{0.664} & \num{0.568} \\
z2m2 & \num{0.854} & \num{0.686} & \num{0.750} & \num{0.713} & \num{0.774} & \num{0.678} & \num{0.514} \\
z2m3 & \num{0.675} & \num{0.720} & \num{0.584} & \num{0.822} & \num{0.748} & \num{0.681} & \num{0.516} \\
z2m4 & \num{0.532} & \num{0.954} & \num{0.614} & \num{0.719} & \num{0.617} & \num{0.637} & \num{0.543} \\
z2m5 & \num{0.406} & \num{0.706} & \num{0.760} & \num{0.750} & \num{0.615} & \num{0.686} & \num{0.580} \\
z2m6 & \num{0.781} & \num{0.875} & \num{0.726} & \num{0.735} & \num{0.726} & \num{0.643} & \num{0.549} \\
\bottomrule

\end{tabular}
\end{table}
Zgodnie z konspektem chłonność wyznaczono ze wzoru Sabine'a,
z przyjętym w instrukcji współczynnikiem 0,16:
\begin{equation}
A(f)=\frac{0{,}16V}{\overline T(f)},\qquad
K_2(f)=10\lg\left(1+\frac{4S}{A(f)}\right).
\end{equation}
Najpierw uśredniono czasy, a następnie przeliczono chłonność.
\begin{table}[H]
\centering\small
\caption{Średnie T20, chłonność i poprawka środowiskowa dla powierzchni na środku pomieszczenia.}
\begin{tabular}{@{}r r r r r r@{}}
\toprule
$f$ (Hz) & $\overline T_{M1}$ (s) & $\overline T_{M2}$ (s) & $\overline T$ (s) & $A$ (m$^2$) & $K_2$ (dB)\\
\midrule
\num{125} & \num{0.544} & \num{0.659} & \num{0.602} & \num{35.112} & \num{4.643} \\
\num{250} & \num{0.828} & \num{0.760} & \num{0.794} & \num{26.591} & \num{5.473} \\
\num{500} & \num{0.727} & \num{0.679} & \num{0.703} & \num{30.039} & \num{5.100} \\
\num{1000} & \num{0.752} & \num{0.760} & \num{0.756} & \num{27.946} & \num{5.319} \\
\num{2000} & \num{0.682} & \num{0.701} & \num{0.692} & \num{30.542} & \num{5.050} \\
\num{4000} & \num{0.644} & \num{0.665} & \num{0.654} & \num{32.269} & \num{4.887} \\
\num{8000} & \num{0.517} & \num{0.545} & \num{0.531} & \num{39.768} & \num{4.296} \\
\bottomrule

\end{tabular}
\end{table}

\subsection{Porównanie metod i potwierdzenie kwalifikacji}
Dla $K_{2A}$ konspekt nakazuje przyjąć czas pogłosu przy 1 kHz
\autocite{konspektKwalifikacja}. Otrzymano
$\overline T(1\,\mathrm{kHz})=\reverbTime\;\unit{\second}$,
$A=\reverbAbsorption\;\unit{\metre\squared}$ oraz
$K_{2A}=\reverbCorrection\;\unit{\decibel}$ dla źródła na środku.
Średnie T20 z dwóch ustawień mikrofonu wynoszą odpowiednio
\qty{0.752}{\second} i \qty{0.760}{\second}.
Dla wszystkich 12 wyników T20 przy 1 kHz moduł współczynnika korelacji
wynosi co najmniej 0,99. Jest to kontrola liniowości dopasowania,
która nie zastępuje oceny zakresu dynamiki i samego zaniku.
Jako kontrolę wyboru parametru obliczono również średnie T30 przy 1 kHz:
\qty{0.731}{\second}, co dawałoby $K_{2A}\approx\qty{5.22}{\decibel}$
dla środka. Oba warianty prowadzą do tej samej oceny kwalifikacji.
\begin{table}[H]
\centering
\caption{Porównanie poprawek A: metoda przybliżona i pogłosowa (T20 przy 1 kHz).}
\begin{tabular}{@{}l r r r@{}}
\toprule
Położenie & $K_{2A,\alpha}$ (dB) & $K_{2A,T}$ (dB) & Różnica $T-\alpha$ (dB)\\
\midrule
Środek & \num{5.683} & \num{5.319} & \num{-0.364} \\
Ściana & \num{3.894} & \num{3.602} & \num{-0.293} \\
Róg & \num{2.601} & \num{2.380} & \num{-0.221} \\
\bottomrule

\end{tabular}
\end{table}
Metoda pogłosowa daje chłonność około 12,4\% większą niż
\qty{24.87}{\metre\squared} uzyskane z $\alpha=\num{0.15}$.
Odpowiada jej efektywny współczynnik przy 1 kHz
$\alpha_{\mathrm{ef}}=A/S_v\approx\reverbAlpha$.
Nie jest to częstotliwościowo niezależny współczynnik materiałowy.
Różnice poprawek nie przekraczają \qty{0.37}{\decibel}.
Obie metody potwierdzają spełnienie kryterium
$K_{2A}\leq\qty{7}{\decibel}$ dla rozpatrywanych powierzchni,
a więc kwalifikację do zastosowanej metody orientacyjnej.
Nie oznacza to automatycznie kwalifikacji do metody technicznej.

W dalszych obliczeniach mocy pozostawiono pierwotne poprawki z metody
przybliżonej, aby porównanie stanowiło osobne sprawdzenie kwalifikacji.
Zastosowanie poprawek pogłosowych zwiększyłoby oszacowania $L_{WA}$
o około 0,36 dB (środek), 0,29 dB (ściana) i 0,22 dB (róg).
